% Transcription — fonds Alexandre Grothendieck, shelfmark 112, batch 1 (pages 1-4).
% Established from the facsimile archives/batches/112/batch-01.pdf (a mirror of
% https://grothendieck.umontpellier.fr/112.pdf), per the skill
% transcribe-grothendieck.
% This is the hardest folder of its group to read, and the transcription says so
% rather than hiding it. Four leaves of computer listing paper written across in
% a very small, very fast hand, the pale printed columns of the listing showing
% through every line. The skeleton is firm — the title, the two definitions
% (lisse, négligeable), the three numbered questions, the A)/B)/C) headings — and
% it is given. A good deal of what hangs off that skeleton is not firm, and is
% marked rather than guessed: the proportion of \ill here is high, and lowering
% it would mean inventing.
% Pass: Opus 5 (claude-opus-5), 2026-08-16 — first pass, unchecked against the pages by a human.
\documentclass[11pt,a4paper]{article}
% Le préambule commun à toutes les transcriptions du fonds Grothendieck.
%
% Il tient en une idée : **l'incertitude se compose, elle ne se résout pas.**
% Une transcription qui lisse un mot douteux a détruit la seule chose pour
% laquelle on la faisait — un lecteur doit pouvoir savoir, à chaque endroit,
% ce qui était sur le feuillet et ce qui a été ajouté. Les six macros qui
% suivent sont donc l'appareil critique, et non de la décoration.
%
% Les mêmes commandes sont reconnues par `scripts/render.mjs`, qui produit la
% vue de lecture du site. Une macro ajoutée ici doit l'être aussi là-bas,
% faute de quoi le rendu échouera bruyamment — ce qui est voulu : un rendu
% silencieusement faux est pire qu'un rendu absent.

\ProvidesPackage{grothendieck}[2026/08/08 Transcriptions du fonds Grothendieck]

\usepackage{amsmath,amssymb,amsthm}
\usepackage{mathrsfs}
\usepackage{stmaryrd}
% Les diagrammes commutatifs sont partout dans ce fonds : c'est la langue
% même de la géométrie algébrique de Grothendieck, et une transcription qui
% les rendrait en prose ne transcrirait rien.
\usepackage{tikz-cd}
% Français par défaut — c'est la langue des feuillets — mais l'anglais est
% chargé aussi : la traduction et le résumé sont des documents anglais, et la
% typographie française (espaces avant les deux-points) y serait une faute.
% Ces documents basculent par \selectlanguage{english} en tête de corps.
\usepackage[english,french]{babel}
\usepackage{fontspec}
\usepackage{geometry}
\geometry{a4paper,margin=25mm,marginparwidth=28mm}
\usepackage{marginnote}
\usepackage{xcolor}
% ulem, not soul. `soul`'s \st and \ul are built on a character-by-character
% scan that XeLaTeX's font machinery breaks: the document fails with an
% undefined control sequence rather than degrading. ulem does the same two jobs
% and is Unicode-engine safe. `normalem` stops it hijacking \emph, which the
% transcriptions use for Grothendieck's own emphasis.
\usepackage[normalem]{ulem}
\usepackage{hyperref}
\hypersetup{colorlinks=true,linkcolor=black,urlcolor=black,pdfborder={0 0 0}}

% --- Métadonnées du lot -----------------------------------------------------
% Reprises telles quelles par le script de rendu, qui les affiche en tête de
% la vue de lecture. Elles fixent ce que le fichier prétend couvrir : sans
% elles, une transcription flotte sans point d'attache dans un fonds de seize
% mille feuillets.

\newcommand{\folder}[1]{\gdef\grothFolder{#1}}
\newcommand{\batch}[1]{\gdef\grothBatch{#1}}
\newcommand{\leaves}[2]{\gdef\grothFirst{#1}\gdef\grothLast{#2}}
\newcommand{\foldertitle}[1]{\gdef\grothFoldertitle{#1}}
\gdef\grothFolder{?}\gdef\grothBatch{?}\gdef\grothFirst{?}\gdef\grothLast{?}\gdef\grothFoldertitle{}

% --- Le résumé --------------------------------------------------------------
%
% Il ouvre la lecture modernisée, sous le titre « Résumé », et il est composé
% en petit. Ce n'est pas un ornement : c'est le seul passage du document qui
% ne soit pas des mathématiques, et un lecteur doit pouvoir le reconnaître
% avant de l'avoir lu — puis le sauter s'il connaît déjà le sujet, ou s'y
% arrêter s'il ne le connaît pas. Le filet qui le suit marque la reprise du
% corps.

\newenvironment{resume}
  {\small\setlength{\parskip}{0.45em}}
  {\par\medskip\hrule\medskip}

% --- Mots-clés ---------------------------------------------------------------
%
% En anglais, en clôture du résumé de la lecture modernisée : le vocabulaire
% moderne sous lequel chercher ce que les feuillets construisent. Le manifeste
% du site (`npm run manifest`) les extrait pour étiqueter la cote entière —
% cette ligne est la source unique des tags, il n'y a pas de fichier à côté.
\newcommand{\keywords}[1]{\par\smallskip\noindent
  {\footnotesize\textsc{Keywords} --- \textit{#1}}\par}

% --- L'appareil critique ----------------------------------------------------

% Le numéro de feuillet, en marge. C'est l'ancre qui permet de régler un
% désaccord sur une formule en regardant *un* feuillet plutôt qu'une chemise
% de six cents. Le site s'en sert aussi pour tourner les pages du fac-similé
% quand on fait défiler la transcription.
\newcommand{\leaf}[1]{%
  \marginnote{\footnotesize\color{gray!70}\textsf{#1}}%
  \label{leaf:#1}\ignorespaces}

% Illisible : on ne devine pas. Un mot inventé qui se lit comme les autres est
% le pire résultat possible.
\newcommand{\ill}{\textcolor{red!60!black}{[\,\ldots\,]}}

% Lecture douteuse : le mot est donné, et signalé comme incertain.
\newcommand{\uncertain}[1]{\uline{#1}}

% Ajout de l'éditeur — une lettre manquante, un mot sous-entendu.
\newcommand{\add}[1]{\textcolor{blue!50!black}{[#1]}}

% Biffé par Grothendieck lui-même. Ce qu'il a rayé fait partie du document :
% une rature dit souvent où la pensée a changé de direction.
\newcommand{\struck}[1]{\sout{#1}}

% Note du transcripteur, sur l'état matériel du feuillet.
\newcommand{\note}[1]{\par\smallskip{\small\color{gray!60!black}\textit{[#1]}}\par\smallskip}

% Note marginale de Grothendieck, distincte du corps du texte.
\newcommand{\marginal}[1]{\par\smallskip\noindent\hspace{1em}%
  {\small\color{gray!60!black}#1}\par\smallskip}

% --- Titre ------------------------------------------------------------------

\AtBeginDocument{%
  \begin{center}
    {\large\bfseries Fonds Alexandre Grothendieck --- cote n\textsuperscript{o}~\grothFolder}\\[2pt]
    {\itshape\grothFoldertitle}\\[4pt]
    {\small feuillets \grothFirst--\grothLast\ (lot \grothBatch)}\\[2pt]
    {\footnotesize\color{gray!60!black}Transcription établie d'après le fac-similé numérisé
     par l'Université de Montpellier.\\
     Les lectures incertaines sont soulignées, les illisibles notées [\,\ldots\,],
     les ajouts entre crochets.}
  \end{center}
  \vspace{1em}}

\endinput


\folder{112}
\batch{1}
\pages{1}{4}
\dating{s.d.}
\watermark{Édition de démonstration}
\foldertitle{Linéarisation : notes manuscrites (s.d.)}

\begin{document}

\page{1}

\section*{Linéarisation}

\note{cette page ne porte que ce titre, à l'encre en haut à droite ; il est
répété au même endroit de la page suivante, dans une équerre tracée à la plume}

\page{2}

Soit $L \in \mathrm{Ob}\,\hat{A}_{\mathrm{ab}}$.

On dit que $L$ est \underline{lisse} si, pour tout $a \to b$ dans $A$,
\[
\int_{A_{a}} L_{a} \longrightarrow \int_{A_{b}} L_{b}
\]
est un \ill, les deux membres étant respectivement égaux à
\[
\int_{A} L^{(a)} \quad\text{et}\quad \int_{A} L^{(b)},
\qquad\text{soit encore}\qquad
L^{(a)} * L^{A^{\circ}} \quad\text{et}\quad L^{(b)} * L^{A^{\circ}}.
\]

\note{les deux colonnes d'égalités sont écrites l'une sous l'autre sur la page,
reliées par des doubles barres ; on les met ici en ligne. Le mot qui qualifie la
flèche --- isomorphisme, quasi-isomorphisme --- est écrit très vite et n'est pas
assuré}

On dit que $L$ est \underline{négligeable} si $L^{b}$ est \ill

c'est-à-dire si
\[
\int_{A_{a}} L \;=\; \int_{A} L^{(a)} \longrightarrow \ill
\]

Soit donc $L^{b} = \bigl(a \mapsto L^{(a)} \otimes \ill\bigr)$ dans
$D_{\bullet}(\hat{B}_{\mathrm{ab}})$. C'est un complexe \ill\ dans
$\hat{B}_{\mathrm{ab}}$ qui est \ill\ à coefficients \uncertain{tordus}.

Soit \struck{\ill} $\mathrm{Lisse}_{\mathbb{Z}}(A)$ la sous-catégorie pleine de
$\hat{A}_{\mathrm{ab}}$ formée des objets lisses.

\begin{enumerate}
\item[(1)] \struck{Lisse} $W_{y}^{-1}\,\mathrm{Lisse}_{\mathbb{Z}}(A)
\longrightarrow D^{\bullet}_{lc}(\hat{B}_{\mathrm{ab}})$ : est-ce une
\underline{équivalence de catégories} ?

\note{sous la flèche, une seconde ligne porte
$D^{\bullet}_{lc}(\hat{A}_{\mathrm{ab}})$ et, en regard, une question sur la
définition, où se lit encore « comme $M_{\bullet} \mapsto M^{b}$ ». OK pour
$A = \Delta$ ; OK pour $A = \Delta_{/X}$, où $X \in \mathrm{Ob}\,\hat{\Delta}$ ?}

\item[(2)] Si $L, M \in \mathrm{Lisse}_{\mathbb{Z}}(A)$, a-t-on
\[
\mathrm{Ext}^{i}_{\mathbb{Z}_{A}}(L, M) \longrightarrow
\mathrm{Ext}^{i}_{\mathbb{Z}_{B}}(L^{b}, M^{b}) \quad ?
\]

\note{en regard, entre crochets : « serait-ce vrai sans supposer $L$, $M$
lisses ? », et, plus haut, « $D_{\bullet}(\mathrm{Ab}) \cong \mathrm{Hot}_{ab}$
si $A$ est asphérique »}

\item[(3)] Supposons $L$ lisse \ill
\end{enumerate}

\page{3}

\[
L_{\bullet} \longmapsto \varphi(L_{\bullet}), \qquad
H_{i}(L_{\bullet}) \simeq H_{i}\,\uncertain{c}\varphi(L_{\bullet})
\]

\begin{itemize}
\item[A)] $D^{b}_{lc}(\hat{A}_{k})$ etc.

\item[B)] $D^{-}_{lc}(\hat{A}_{h}) \rightleftarrows D^{-}_{lc}(\hat{B}_{h})$
--- \underline{équiv.} \ill

\note{NB : on a $D^{-}(\hat{A}_{k}) \to D^{-}(\hat{B}_{h})$ \ill}

\item[C)] Objets \underline{lisses} de $\hat{A}_{k}$ (plus général) :
\struck{$\mathrm{Lisse}$}$_{k}$, $\hat{A}_{k\text{-lisse}}$, $W$
\end{itemize}

\page{4}

\note{la dernière page écrite reprend les mêmes matières sans les mener plus
loin : on y relit $\mathrm{Lisse}_{\mathbb{Z}}(A)$, la flèche vers
$D_{lc}^{-}$, et le cas $A = \Delta$. Elle est plus raturée que les
précédentes et rien n'y forme un énoncé qu'on puisse détacher ; on ne la
découpe donc pas en propositions qu'elle ne porte pas}

\end{document}
